\section{恒等态射与等价}\label{section:2}
\subsection{定义}\label{sec:2a}
设 $\A$ 是非含幺 $\Ai$ 范畴。“$\A$ 具有恒等态射”有几种不同含义。

\begin{definition*}
若对每个对象 $X$，都有一个必然唯一的 $e_X\in\hom^0_{\A}(X,X)$，满足
\begin{equation}\label{eq:2.1}
\begin{gathered}
\mu^1_{\A}(e_X)=0;\\
(-1)^{|a|}\mu^2_{\A}(e_{X_1},a)=a=\mu^2_{\A}(a,e_{X_0})\qquad\text{对 }a\in\hom_{\A}(X_0,X_1);\\
\mu^d_{\A}(a_{d-1},\ldots,a_{n+1},e_{X_n},a_n,\ldots,a_1)=0\\
\text{对 }d>2,\ a_k\in\hom_{\A}(X_{k-1},X_k)\text{ 及任意 }0\le n<d,
\end{gathered}
\end{equation}
则称 $\A$ \emph{严格含幺}。
\end{definition*}

\begin{definition*}
若 $H(\A)$ 含幺，即对每个对象 $X$ 都有恒等态射，从而是通常意义下的分次线性范畴，则称 $\A$ \emph{上同调含幺}，或简称 \emph{c-含幺}。
\end{definition*}

\SourcePage{31}
\begin{definition*}
\emph{同伦含幺 $\Ai$ 范畴}由一个集合 $\Ob\A$、分次向量空间 $\hom_{\A}(X_0,X_1)$，以及对所有 $d+i_0+\cdots+i_d>0$ 给出的多重线性映射组成：
\begin{equation}\label{eq:2.2}
\begin{aligned}
\mu^{d,(i_d,\ldots,i_0)}_{\A}:\hom_{\A}(X_{d-1},X_d)\otimes\cdots\otimes\hom_{\A}(X_0,X_1)\\
\longrightarrow\hom_{\A}(X_0,X_d)\bigl[2-d-2\textstyle\sum_k i_k\bigr].
\end{aligned}
\end{equation}
这些映射须满足某些推广的结合律方程；当 $i_0=\cdots=i_d=0$ 时，它们归结为\eref{1.2}。最先出现的几个新方程是
\begin{equation}\label{eq:2.3}
\begin{gathered}
\mu^1_{\A}(\mu^{0,(1)}_{\A})=0,\\
(-1)^{|a|-1}\mu^2_{\A}(\mu^{0,(1)}_{\A},a)
+\mu^1_{\A}(\mu^{1,(1,0)}_{\A}(a))
+\mu^{1,(1,0)}_{\A}(\mu^1_{\A}(a))=a,\\
\mu^2_{\A}(a,\mu^{0,(1)}_{\A})
+\mu^1_{\A}(\mu^{1,(0,1)}_{\A}(a))
+(-1)^{|a|-1}\mu^{1,(0,1)}_{\A}(\mu^1_{\A}(a))=-a,\\
\mu^{1,(1,0)}_{\A}(\mu^{0,(1)}_{\A})
+\mu^{1,(0,1)}_{\A}(\mu^{0,(1)}_{\A})
+\mu^1_{\A}(\mu^{0,(2)}_{\A})=0.
\end{gathered}
\end{equation}

\end{definition*}
除去通常的符号因素，前面的两条乘法关系说明，与上闭链 $e_X=-\mu^{0,(1)}_{\A}$ 作左乘或右乘，都链同伦于恒等映射。利用这两个同伦中的任意一个，可见 $\mu^2_{\A}(e_X,e_X)$ 与 $e_X$ 只差一个上边界。这两个上边界之差是上闭链，而\eref{2.3}的最后一部分表明，它本身在上同调中也是平凡的。完整的方程组可理解为把这一层级继续推进，依次给出越来越高阶的同伦；细节见\bcite{55}，第 5 节。

任意严格含幺 $\Ai$ 范畴都是同伦含幺的：令 $\mu^{0,(1)}_{\A}=-e_X$，并把\eref{2.2}中所有其他满足 $i_0+\cdots+i_d>0$ 的映射取为零即可。此外，如前面的讨论所示，同伦含幺蕴含 c-含幺。\znote{23}{这三种“含幺”其实处在不同层次：严格含幺要求单位元在链层面就满足结合律，同伦含幺只要求它以链同伦的方式充当单位元，而上同调含幺只要求 $H(\A)$ 含有单位元。几何中自然出现的 Fukaya 范畴通常只是上同调含幺的，故后面以它为默认框架，必要时再用形式微分同胚换成严格含幺的模型。[来源：译者按（上下文推断，无外部依据）：依据本页 (2.1)–(2.3) 及下页对三种概念差别的讨论]}

\begin{definition*}
还应提到函子的相应概念。若 $\A,\B$ 严格含幺，则称 $\Ai$ 函子 $\F:\A\to\B$ \emph{严格含幺}，如果对所有对象 $X$ 都有 $\F^1(e_X)=e_{\F(X)}$，并且对任意 $d\ge2$ 及任意 $n$，都有
\[
\F^d(a_{d-1},\ldots,a_{n+1},e_{X_n},a_n,\ldots,a_1)=0.
\]
若 $\A,\B$ 为 c-含幺，则称 $\Ai$ 函子 $\F:\A\to\B$ 为 \emph{c-含幺}，只需 $H(\F)$ 含幺。同伦含幺版本见\bcite{55}。
\end{definition*}

严格含幺在许多方面很方便，因此在文献中广泛使用，尽管它看起来可能与同伦代数结构的一般理念不太协调。c-含幺乍看也可能令人怀疑，因为它走向另一个极端，把问题完全推到上同调层次。就我们预期的应用而言，Fukaya 范畴并不自然带有严格单位元，但很容易看出它是 c-含幺的。事实上，可以通过几何方式为它构造同伦单位元，利用附加的模空间定义\eref{2.2}中所需的、满足 $i_0+\cdots+i_d>0$ 的映射，见\bcite{60}，§20。所有这些似乎都表明同伦含幺是最恰当的概念，然而下面的讨论将出现一个或许出人意料的转折：
\SourcePage{32}
考察相关障碍后发现，存在一般的代数步骤，可以给任意 c-含幺 $\Ai$ 范畴配备同伦单位元。此外，由\bcite{55}中的推导可见，任意同伦含幺 $\Ai$ 范畴都与一个严格含幺的范畴拟同构。这意味着，三种概念之间的区别在于偏好和便利，而非实质。我们选择主要在 c-含幺的框架中工作。不过，作为证明一些基本性质的技术工具，仍将给出把 c-含幺 $\Ai$ 范畴变为严格含幺范畴的直接方法。\znote{24}{把只在上同调层面成立的单位元提升为链层面的同伦单位元，是同调扰动理论的典型应用：先在同调上选好代表，再逐阶修正高阶映射以抵消障碍，最后由收敛的无穷级数得到新结构。这里的形式微分同胚正是这种“扰动”，引理 2.1 的递归构造则是它的具体实现。[来源：nLab "homological perturbation theory"，https:/\allowbreak /\allowbreak ncatlab.org/\allowbreak nlab/\allowbreak show/\allowbreak homological+\allowbreak perturbation+\allowbreak theory（2026-11-15 访问）]}

\begin{lemma}\label{st:2.1}
设 $\A$ 是上同调含幺 $\Ai$ 范畴，并且对每个对象 $X$，或者复形 $\hom_{\A}(X,X)$ 具有非零上同调，或者对所有 $Y$ 都有 $\hom_{\A}(X,Y)=\hom_{\A}(Y,X)=0$。那么存在满足 $\Phi^1=\Id$ 的形式微分同胚 $\Phi$，使修改后的 $\Ai$ 结构 $\Phi_*\A$ 严格含幺。\znote{24}{原书引理 2.1 没有第二个条件，但该条件不可省：若范畴含零对象 $X$，则 $\hom_{\A}(X,X)$ 上同调为零，代表恒等态射的闭链 $e_X$ 只能是 $0$，而严格含幺要求 $\mu^2(e_X,a)=a$、$\mu^2(a,e_X)=a$，矛盾；故原陈述在含零对象时失效。译文补上的条件正是为排除这一退化情形：当 $X$ 与一切对象的双向 hom 复形都为零时，$e_X=0$ 已使严格含幺方程空洞地成立。此系对原书陈述的必要修补，不是原文内容。[来源：原书 p.~24 引理 2.1；反例由本书 \S1 的 $\Ai$ 范畴定义直接验证，无外部依据。]}
\end{lemma}
\begin{proof}
第一步如下。对每个对象 $X$，选择一个上闭链 $e_X\in\hom^0_{\A}(X,X)$，代表 $H(\A)$ 中的恒等态射。对于任意三个对象，可以找到一个链映射
\[
\mu^2_{\widetilde{\A}}:\hom_{\A}(X_1,X_2)\otimes\hom_{\A}(X_0,X_1)\longrightarrow\hom_{\A}(X_0,X_2),
\]
使它代表 $H(\A)$ 中的复合，并满足\eref{2.1}的第二个条件。取一个形式微分同胚，令 $\Phi^1=\Id$，并令每个 $\Phi^2$ 都是 $\mu^2_{\A}$ 与 $\mu^2_{\widetilde{\A}}$ 之间的链同伦。那么 $\widetilde{\A}=\Phi_*\A$ 的前两个复合映射就是 $\mu^1_{\widetilde{\A}}=\mu^1_{\A}$ 和选定的 $\mu^2_{\widetilde{\A}}$。

从这里起，构造递归进行。设非含幺 $\Ai$ 范畴 $\A$ 带有元素 $e_X\in\hom^0_{\A}(X,X)$，满足\eref{2.1}的前两个等式。考虑对这样的 $\A$ 给出的一族条件 $(U_{d,n})$，其中 $d>2$、$0\le n\le d$，且只考察阶数 $i>2$：
\[
\begin{aligned}
(U_{d,n})\quad&\mu^i_{\A}(a_{i-1},\ldots,a_{j+1},e_{X_j},a_j,\ldots,a_1)\text{ 在以下两种情形为零：}\\
&i<d\text{ 且 }j\text{ 任意；以及 }i=d\text{ 且 }j<n.
\end{aligned}
\]
纳入 $n=0$ 的情形纯粹是为了记号方便，因为显然 $(U_{d,0})=(U_{d-1,d-1})$。设 $\A$ 已对某个 $n<d$ 满足 $(U_{d,n})$。取一个形式微分同胚 $\Phi$，使 $\Phi^1=\Id$，对 $2\le k<d-1$ 有 $\Phi^k=0$，并令
\begin{equation}\label{eq:2.4}
\begin{aligned}
\Phi^{d-1}(a_{d-1},\ldots,a_1)&=(-1)^{\sgnmark_n}\mu^d_{\A}(a_{d-1},\ldots,a_{n+1},e_{X_n},a_n,\ldots,a_1),\\
\Phi^d(a_d,\ldots,a_1)&=(-1)^{\sgnmark_n}\mu^{d+1}_{\A}(a_d,\ldots,a_{n+1},e_{X_n},a_n,\ldots,a_1).
\end{aligned}
\end{equation}
由定义显然可见，在低于 $d-1$ 阶时，$\widetilde{\A}=\Phi_*\A$ 的复合映射与 $\A$ 的相同。利用 $\Ai$ 结合律方程及归纳假设，作一个略复杂的计算，可知 $d-1$ 阶时也相同。现在考察 $d$ 阶；为简化公式，记
\SourcePage{33}
$(b_1,\ldots,b_{d+1})=(a_1,\ldots,a_n,e_{X_n},a_{n+1},\ldots,a_d)$。可得
\begin{equation}\label{eq:2.5}
\begin{aligned}
&\mu^d_{\widetilde{\A}}(a_d,\ldots,a_1)\\
&=\mu^d_{\A}(a_d,\ldots,a_1)
-(-1)^{\S}\mu^2_{\A}\bigl(\mu^d_{\A}(a_d,\ldots,e_{X_{n+1}},a_{n+1},\ldots,a_2),a_1\bigr)\\
&\quad-(-1)^{\sgnmark_{n+1}}\sum_{j<n}(-1)^{\dagger}
\mu^d_{\A}\bigl(a_d,\ldots,e_{X_{n+1}},a_{n+1},\ldots,\mu^2_{\A}(a_{j+2},a_{j+1}),\ldots,a_1\bigr)\\
&\quad-(-1)^{\sgnmark_n}\sum_{i,j}(-1)^{\ddagger}
\mu^{d+2-i}_{\A}\bigl(b_{d+1},\ldots,\mu^i_{\A}(b_{j+i},\ldots,b_{j+1}),b_j,\ldots,b_1\bigr),
\end{aligned}
\end{equation}
其中符号为 $\S=|a_2|+\cdots+|a_{n+1}|-n$、$\dagger=|a_1|+\cdots+|a_j|-j$、$\ddagger=|b_1|+\cdots+|b_j|-j$；最后一个和式取遍以下四种可能：$i=1$ 且 $j\ne n$；$i=2$ 且 $j>n$；$i=d$ 且 $j=0$；最后是 $i=d+1$。注意，按归纳假设，最后一个和式中的表达式在 $i=1$、$j=n$ 时，以及 $3\le i\le d-1$ 时，也都为零。

先设 $n=0$，此时\eref{2.5}右侧第三个表达式为空和。去掉它，并对第四个表达式应用 $\Ai$ 结合律方程，最终得到
\[
\mu^d_{\widetilde{\A}}(a_d,\ldots,a_1)
=\mu^d_{\A}(a_d,\ldots,a_1)-\mu^2_{\A}\bigl(\mu^d_{\A}(a_d,\ldots,a_2,e_{X_1}),a_1\bigr).
\]
当 $a_1$ 也是恒等态射时，此式显然为零，表明 $\widetilde{\A}$ 具有性质 $(U_{d,1})$。接着考虑 $n>0$ 的情形。此时，\eref{2.5}最后一个表达式中，令 $i=2$ 并对所有 $j\le n$ 求和所得的和也为零；$i=d$、$j=1$ 时的项同样为零。像前面一样利用 $\Ai$ 结合律方程，可证明整个和式为零。消去这一和式之后，很容易验证 $\widetilde{\A}$ 具有性质 $(U_{d,n+1})$。

反复进行这一过程，得到一列形式微分同胚 $\Phi_1,\Phi_2,\ldots$，使得当 $k\to\infty$ 时，$(\Phi_k\circ\cdots\circ\Phi_1)_*\A$ 满足越来越多的严格含幺条件。$\Phi_k$ 本身越来越接近恒等映射；也就是说，当 $k\to\infty$ 时，$\Phi_k$ 的第一个非平凡项的阶数越来越高。因此，无穷复合 $\Phi=\cdots\circ\Phi_3\circ\Phi_2\circ\Phi_1$ 是有意义的，且 $\Phi_*\A$ 严格含幺，正如所要证明的那样。\znote{25}{形式微分同胚保持对象不变，只逐阶修改高阶复合映射，因此它是 $\A$ 与 $\Phi_*\A$ 之间在对象上为恒等的 $\Ai$ 同构。证明的每一轮都把第一个非平凡项的阶数推高，故无穷复合在每一阶上只有有限多个非零项，收敛性由此得到保证。[来源：nLab "homological perturbation theory"，https:/\allowbreak /\allowbreak ncatlab.org/\allowbreak nlab/\allowbreak show/\allowbreak homological+\allowbreak perturbation+\allowbreak theory（2026-11-15 访问）]}
\end{proof}

\begin{note}\label{st:2.2}
\stref{引理}{2.1}由 Lefèvre\bcite{95}，定理 3.2.1.1，与作者各自独立得到。还有一个更具概念性的证明，基于同调扰动理论，以及含幺代数的完整 Hochschild 复形与约化 Hochschild 复形拟同构这一事实（表达式\eref{2.4}受到这些复形之间经典同伦等价的启发）。$\Ai$ 函子、$\Ai$ 模等也有类似结果，分别见\bcite{95}，定理 3.2.2.1 和定理 3.3.1.2。
\end{note}

\SourcePage{34}
\subsection{函子范畴}\label{sec:2b}
若 $\A$ 是严格含幺 $\Ai$ 范畴，那么对任意非含幺 $\Ai$ 范畴 $\C$，$\nufun(\C,\A)$ 仍然严格含幺。具体地，严格恒等态射是自然变换 $E_{\F}$，定义为 $E^0_{\F}=e_{\F(X)}\in\hom^0_{\A}(\F(X),\F(X))$，以及对所有 $d>0$ 令 $E^d_{\F}=0$。还有一个稍不显然的结论。

\begin{lemma}\label{st:2.3}
若 $\A$ 为 c-含幺，则 $\Q=\nufun(\C,\A)$ 也为 c-含幺。此外，若 $E_{\F}\in\hom^0_{\Q}(\F,\F)$ 是代表 $H(\Q)$ 中恒等态射的任意链，那么对每个 $X\in\Ob\C$，$E^0_{\F}\in\hom^0_{\A}(\F(X),\F(X))$ 代表 $H(\A)$ 中的恒等态射。
\end{lemma}
\begin{proof}
可将其归结为严格含幺的情形。先利用\stref{笔记}{1.13}取 $\A$ 的最小模型，再应用\stref{引理}{2.1}严格化，得到与 $\A$ 拟同构的严格含幺范畴 $\widetilde{\A}$，以及一个在对象上为恒等的拟同构 $\Phi:\A\to\widetilde{\A}$。通过诱导的上同调同构，将 $H(\widetilde{\A})$ 与 $H(\A)$ 识别。令 $\widetilde{\Q}=\nufun(\C,\widetilde{\A})$。考虑与 $\Phi$ 作左复合的 $\Ai$ 函子，并记 $\widetilde{\F}=\Phi\circ\F$。在自然变换层次上，得到交换图
\SourceEquationDiagram{diagram-2-6}{0.55038403}{2.6}
这里 $C=H(\C)$、$A=H(\A)=H(\widetilde{\A})$、$F=H(\F)=H(\widetilde{\F})$ 是上同调层次的结构，向下的箭头则是\eref{1.10}中的映射，把 $\Ai$ 函子的自然变换送到相应的上同调自然变换。\stref{引理}{1.7}说明 $H(\LG^1_{\Phi})$ 是同构。选择一个链 $E_{\F}$，使 $[\LG^1_{\Phi}(E_{\F})]$ 为严格单位元的类 $[E_{\widetilde{\F}}]$。由 $H(\LG^1_{\Phi})$ 的函子性，$[E_{\F}]$ 是 $H(\Q)$ 中的恒等态射。在\eref{2.6}中追图可知，$H(E_{\widetilde{\F}})=H(E_{\F})$ 是 $\operatorname{Nu\text{-}Fun}(C,A)$ 中的恒等自然变换。这就表明，对每个 $X$，在上述识别下 $[E^0_{\widetilde{\F}}]=[E^0_{\F}]$ 是 $A$ 中的恒等态射，正如所要证明的那样。
\end{proof}

\begin{lemma}\label{st:2.4}
设 $\G:\A\to\B$ 是 c-含幺 $\Ai$ 函子。那么对任意非含幺 $\Ai$ 范畴 $\C$，复合函子 $\LG_{\G}:\nufun(\C,\A)\to\nufun(\C,\B)$ 为 c-含幺。（右复合也有类似结果：设 $\G:\A\to\B$ 是非含幺 $\Ai$ 范畴之间的非含幺 $\Ai$ 函子。那么，对任意 c-含幺 $\Ai$ 范畴 $\C$，复合函子 $\RG_{\G}:\nufun(\B,\C)\to\nufun(\A,\C)$ 为 c-含幺。）
\end{lemma}
\begin{proof}
设 $E_{\F}\in\hom^0_{\nufun(\C,\A)}(\F,\F)$ 是代表 $H(\nufun(\C,\A))$ 中恒等态射的链。取其在复合下的像：
\begin{equation}\label{eq:2.7}
[\LG^1_{\G}(E_{\F})]\in\Hom^0_{H(\nufun(\C,\B))}(\G\circ\F,\G\circ\F).
\end{equation}
我们已经看到，对每个 $X$，$E^0_{\F}\in\hom^0_{\A}(\F(X),\F(X))$ 代表上同调恒等态射，因此 $\LG^1_{\G}(E_{\F})^0=\G^1(E^0_{\F})\in\hom^0_{\B}(\G(\F(X)),\G(\F(X)))$ 也代表上同调恒等态射。
\SourcePage{35}
利用\stref{引理}{1.6}及 $\nufun(\C,\B)$ 的 c-含幺性，可知\eref{2.7}是可逆态射。由于 $\LG_{\G}$ 的函子性，它还是幂等的，故必等于恒等态射（后面会反复使用这一类论证）。对于 $\RG_{\G}$，证明相同。
\end{proof}

\subsection{同伦与同构}\label{sec:2c}
设 $\A$ 是严格含幺 $\Ai$ 范畴，$\C$ 是非含幺 $\Ai$ 范畴，而 $\F_0,\F_1:\C\to\A$ 是在对象上作用相同的非含幺 $\Ai$ 函子。设 $T$ 是从 $\F_0$ 到 $\F_1$ 的零次预自然变换，满足 $T^0=0$。用 $S$ 表示修改后的预自然变换，定义为对所有 $X$ 令 $S^0=e_{\F_0(X)}$，且当 $d>0$ 时令 $S^d=T^d$。直接计算表明，$T$ 是 $\F_0$ 与 $\F_1$ 之间的同伦，当且仅当 $S$ 是真正的自然变换。在这一背景下，原本看似神秘的同伦复合表达式\eref{1.16}，可以看作相关自然变换的通常复合；\eref{1.17}也有类似解释。

利用\stref{引理}{1.6}及 $\Ai$ 函子范畴的严格含幺性，可知两个同伦函子作为 $H^0(\nufun(\C,\A))$ 的对象是同构的。我们还需要 c-含幺背景下的类似结果。

\begin{lemma}\label{st:2.5}
设 $\A$ 是 c-含幺 $\Ai$ 范畴，$\C$ 是非含幺 $\Ai$ 范畴。若两个非含幺 $\Ai$ 函子 $\F_0,\F_1:\C\to\A$ 同伦，那么它们作为 $H^0(\nufun(\C,\A))$ 的对象也是同构的。
\end{lemma}
\begin{proof}
沿用\stref{引理}{2.3}证明中的记号。由于左复合保持同伦，$\widetilde{\F}_k=\Phi\circ\F_k$ 同伦，因此由前面的观察，它们作为 $H^0(\widetilde{\Q})$ 的对象同构。另一方面，由\stref{引理}{1.7}，$\LG_{\Phi}$ 上同调全忠实，所以 $\F_k$ 本身必然同构。
\end{proof}

\subsection{Morita 不变性}\label{sec:2d}
设 $A,B$ 是线性分次含幺范畴，$F_0,F_1:A\to B$ 是线性分次含幺函子。设子范畴 $\widetilde A\subset A$ 的包含函子为等价。

\begin{lemma}\label{st:2.6}
限制映射给出 Hochschild 上同调群的同构 $HH(A,B)\cong HH(\widetilde A,B)$。
\end{lemma}
\begin{proof}
设 $\rho:CC(A,B)\to CC(\widetilde A,B)$ 是链层次的限制映射。对任意 $X\in\Ob A$，选择一个 $\widetilde X\in\Ob\widetilde A$，以及同构 $u_X:X\to\widetilde X$。对于已属于 $\widetilde A$ 的对象 $X$，可以取 $\widetilde X=X$，并令 $u_X$ 为恒等态射。\emph{Dennis 余迹}是如下链映射：
\[
\begin{aligned}
\chi &:CC^{r+s}(\widetilde A,B)^s\longrightarrow CC^{r+s}(A,B)^s,\\
(\chi T)^r(a_r,\ldots,a_1)
&=F_1(u_{X_r})^{-1}T^r(u_{X_r}a_ru_{X_{r-1}}^{-1},\ldots,u_{X_1}a_1u_{X_0}^{-1})F_0(u_{X_0}).
\end{aligned}
\]
\SourcePage{36}
它是同伦等价，因为 $\rho\circ\chi=\id$，且 $\id-\chi\circ\rho=\partial\eta+\eta\partial$，其中
\[
\begin{aligned}
(\eta T)(a_{r-1},\ldots,a_1)
&=\sum_n(-1)^{\sgnmark_n+(r+s)}F_1(u_{X_{r-1}})^{-1}\\
&\qquad\cdot T(u_{X_{r-1}}a_{r-1}u_{X_{r-2}}^{-1},\ldots,u_{X_{n+1}}a_{n+1}u_{X_n}^{-1},u_{X_n},a_n,\ldots,a_1).
\end{aligned}
\]
\end{proof}

\begin{note}\label{st:2.7}
也可以采用以下更抽象的途径。考虑分次 $A$ 双模组成的Abel 范畴 $\operatorname{Bimod}^0(A)$；更精确地说，其对象是从 $(A^{\mathrm{opp}},A)$ 到分次向量空间范畴的双函子（双模同态就是零次自然变换）。任意一对函子 $(G_0,G_1)$ 都给出一个双模，由向量空间 $\Hom_B(G_0(X_0),G_1(X_1))$ 组成。略微滥用记号，把这个双模直接记作 $B$，而把考察函子对 $(\Id_A,\Id_A)$ 得到的相应对象记作 $A$。Hochschild 复形可以解释为来自最后这个双模在 $\operatorname{Bimod}^0(A)$ 中的一个投射分解，因此
\begin{equation}\label{eq:2.8}
HH(A,B)\cong\operatorname{Ext}_{\operatorname{Bimod}^0(A)}(A,B).
\end{equation}
这里考察所有阶数的 Ext 群，并把 $B$ 的给定分次作所有可能的平移，所以右侧是双分次的。由初等理由，将双函子从 $A$ 限制到 $\widetilde A$，给出范畴等价 $\operatorname{Bimod}^0(A)\cong\operatorname{Bimod}^0(\widetilde A)$，从而得到 Ext 群的同构。

现在，这与代数的 Hochschild 同调和上同调的经典 Morita 不变性之间的关系应当已经清楚，例如参看\bcite{143}，定理 9.5.6。事实上，我们的情形类似于矩阵代数给出的 Morita 不变性的基本例子，上面的初等论证沿用了\bcite{96}，定理 1.2.4 的证明。
\end{note}

\subsection{拟等价}\label{sec:2e}
\begin{definition*}
还需要补充一点记号和术语：给定两个 c-含幺 $\Ai$ 范畴 $\A,\B$，用 $\operatorname{fun}(\A,\B)$ 表示 $\nufun(\A,\B)$ 中由 c-含幺 $\Ai$ 函子组成的全子范畴。若 c-含幺 $\Ai$ 函子 $\F:\A\to\B$ 的基础上同调函子为等价，则称 $\F$ 为\emph{拟等价}。
\end{definition*}

设 $\A$ 是 c-含幺 $\Ai$ 范畴，$\widetilde{\A}\subset\A$ 为全 $\Ai$ 子范畴，且包含函子为拟等价。
\begin{lemma}\label{st:2.8}
对于任意 c-含幺 $\Ai$ 范畴 $\B$，将 c-含幺 $\Ai$ 函子及其预自然变换作限制，给出拟等价 $\Q=\operatorname{fun}(\A,\B)\to\widetilde{\Q}=\operatorname{fun}(\widetilde{\A},\B)$。
\end{lemma}
\begin{proof}
考虑任意两个 c-含幺 $\Ai$ 函子 $\G_k$（$k=0,1$），以及它们在 $\widetilde{\A}$ 上的限制 $\widetilde{\G}_k$。预自然变换上的限制映射
\begin{equation}\label{eq:2.9}
\hom_{\Q}(\G_0,\G_1)\longrightarrow\hom_{\widetilde{\Q}}(\widetilde{\G}_0,\widetilde{\G}_1)
\end{equation}
\SourcePage{37}
显然与长度滤过相容，因此诱导相应谱序列之间的映射。在 $E_2$ 层次上，它归结为相关上同调函子的 Hochschild 上同调上的限制映射；我们刚刚已经证明，这个映射是同构。因此\eref{2.9}是拟同构，也就是说，限制函子上同调全忠实。最后，给定任意 c-含幺 $\Ai$ 函子 $\widetilde{\G}:\widetilde{\A}\to\B$，可以把其基础上同调函子扩张为 $G:H(\A)\to H(\B)$，并且这一扩张在自然同构意义下唯一；随后\stref{引理}{1.10}允许构造一个扩张 $\widetilde{\G}$、满足 $H(\G)=G$ 的 $\G:\A\to\B$，因而它为 c-含幺。这证明限制函子在对象上为满射，故为拟等价。
\end{proof}

特别地，扩张步骤说明存在 c-含幺 $\Ai$ 函子 $\mathcal P:\A\to\widetilde{\A}$，满足 $\mathcal P|\widetilde{\A}=\Id_{\widetilde{\A}}$。若用 $\mathcal K:\widetilde{\A}\to\A$ 表示包含函子，则 $\mathcal P\circ\mathcal K=\Id_{\widetilde{\A}}$。$\mathcal P$ 是拟等价，特别地上同调全忠实，因此由\stref{引理}{1.7}，$\LG_{\mathcal P}$ 也上同调全忠实。由于 $\LG_{\mathcal P}(\mathcal K\circ\mathcal P)=\LG_{\mathcal P}(\Id_{\A})$，可知 $\mathcal K\circ\mathcal P$ 与 $\Id_{\A}$ 作为 $H^0(\operatorname{fun}(\A,\A))$ 的对象同构。\znote{29}{在链复形层面，诱导同调同构的链映射称为拟同构；它一般不可逆，“存在从 $C$ 到 $D$ 的拟同构”甚至不是对称关系。拟等价要求上同调函子本身为等价，因而天然可逆，\stref{定理}{2.9}正是这一点的定量表述。[来源：nLab "quasi-isomorphism"，https:/\allowbreak /\allowbreak ncatlab.org/\allowbreak nlab/\allowbreak show/\allowbreak quasi-isomorphism（2026-11-15 访问）]}

\begin{theorem}\label{st:2.9}
设 $\F:\A\to\B$ 为拟等价。那么存在拟等价 $\G:\B\to\A$，使得 $\G\circ\F\cong\Id_{\A}$ 在 $H^0(\operatorname{fun}(\A,\A))$ 中成立，且 $\F\circ\G\cong\Id_{\B}$ 在 $H^0(\operatorname{fun}(\B,\B))$ 中成立。
\end{theorem}
\begin{proof}
存在全 $\Ai$ 子范畴 $\widetilde{\A}\subset\A$、$\widetilde{\B}\subset\B$，使包含函子为拟等价，并使 $\F$ 限制为从 $\widetilde{\A}$ 到 $\widetilde{\B}$ 的拟同构 $\widetilde{\F}$。如前所述，两个包含函子都能在各自的函子范畴中取到同构意义下的逆。由\stref{推论}{1.14}，$\widetilde{\F}$ 本身具有同伦意义下的逆；再由\stref{引理}{2.5}，同伦蕴含函子范畴中的同构。
\end{proof}

\begin{lemma}\label{st:2.10}
若 $\G:\A\to\B$ 为拟等价，那么对于任意 c-含幺 $\Ai$ 范畴 $\C$，$\RG_{\G}:\operatorname{fun}(\B,\C)\to\operatorname{fun}(\A,\C)$ 也为拟等价。（$\LG_{\G}$ 也有相应结论；事实上，后一情形还存在适用于非含幺 $\C$ 及范畴 $\nufun$ 的版本。）
\end{lemma}
\begin{proof}
首先需要知道 $\RG_{\G}$ 上同调全忠实。若 $\G$ 是拟同构，这是\stref{引理}{1.7}的一部分；若 $\G$ 是全子范畴的包含函子，这是\stref{引理}{2.8}的一部分。与前面的讨论一样，这两个特殊情形推出一般结论。接着，\stref{定理}{2.9}给出一个 $\F:\B\to\A$，使 $\F\circ\G\cong\Id_{\A}$。由\stref{引理}{2.4}，对任意 $\mathcal H\in\Ob\operatorname{fun}(\A,\C)$，有
\[
\RG_{\G}(\mathcal H\circ\F)=\LG_{\mathcal H}(\F\circ\G)
\cong\LG_{\mathcal H}(\Id_{\A})=\mathcal H,
\]
这表明 $\RG_{\G}$ 在对象的同构类上为满射。$\LG_{\G}$ 的论证类似，但略为简单。
\end{proof}

\begin{note}\label{st:2.11}
\stref{定理}{2.9}的严格含幺版本，似乎最早由 Kontsevich 提出（未发表）。它出现于\bcite{55}，定理 8.6，以及\bcite{95}，定理 9.2.04。这里讨论的 c-含幺版本也见于\bcite{97}，定理 8.8。
\end{note}

\SourcePage{38}
\subsection{\texorpdfstring{$A_\infty$}{A∞} 模}\label{sec:2f}
设 $\A$ 是 c-含幺 $\Ai$ 范畴。由于 $\Ch$ 严格含幺，$\Q=\operatorname{nu\text{-}mod}(\A)=\nufun(\A^{\mathrm{opp}},\Ch)$ 也严格含幺。注意，由于\eref{1.20}中引入的符号，$\M$ 的严格恒等态射是如下模同态 $t=e_{\M}$：
\begin{equation}\label{eq:2.10}
t^1(b)=(-1)^{|b|}b,\qquad t^d=0\quad\text{对所有 }d\geq2.
\end{equation}
\begin{definition*}
若 $\Ai$ 模 $\M$ 的基础上同调模含幺，就称 $\M$ 为 \emph{c-含幺}。这意味着：若 $e_X\in\hom^0_{\A}(X,X)$ 代表上同调单位元，则链映射 $\mu^2_{\M}(\cdot,e_X):\M(X)\to\M(X)$ 在上同调上诱导恒等映射。用 $\modcat(\A)$ 表示由 c-含幺 $\Ai$ 模组成的全子范畴。
\end{definition*}
若 $\G:\A\to\B$ 为 c-含幺 $\Ai$ 函子，可以把拉回\eref{1.22}限制到这些子范畴，得到
\begin{equation}\label{eq:2.11}
\G^*:\modcat(\B)\longrightarrow\modcat(\A),
\end{equation}
它是严格含幺 dg 函子。作为\stref{引理}{2.10}的特殊情形，若 $\G$ 为拟等价，那么\eref{2.11}也为拟等价。

\subsection{Yoneda 嵌入}\label{sec:2g}
由定义，对 c-含幺的 $\A$，Yoneda 嵌入 $\mathcal I=\mathcal I_{\A}$ 的像落在子范畴 $\modcat(\A)\subset\operatorname{nu\text{-}mod}(\A)$ 中。
\begin{lemma}\label{st:2.12}
对任意 c-含幺 $\Ai$ 模 $\M$ 及 $\A$ 的任意对象 $Y$，\eref{1.25}中的映射 $\lambda$ 是拟同构。
\end{lemma}
\begin{proof}
由定义，$\lambda$ 的映射锥平移 $-1$ 后，是链复形
\begin{equation}\label{eq:2.12}
\left(\M(Y)\oplus\hom_{\Q}(\Y,\M)[-1],\begin{pmatrix}\mu^1_{\M}&0\\\lambda&-\mu^1_{\Q}\end{pmatrix}\right).
\end{equation}
先取子复形 $\hom_{\Q}(\Y,\M)[-1]$，再接上它通常的长度滤过，以此给这个复形赋予滤过。记 $A=H(\A)$、$M=H(\M)$，所得谱序列满足
\[
E_1^{rs}=\begin{cases}
M^s(Y),&r=0,\\[3pt]
\displaystyle\prod_{X_0,\ldots,X_{r-1}}\Hom^s_{\K}\!\left(\begin{aligned}&\Hom_A(X_{r-1},Y)\otimes\Hom_A(X_{r-2},X_{r-1})\\[-2pt]&\qquad\otimes\cdots\otimes\Hom_A(X_0,X_1),\,M(X_0)\end{aligned}\right),&r>0.
\end{cases}
\]
由\eref{1.25}和\eref{1.21}可见，微分 $\partial=\partial^{rs}:E_1^{rs}\to E_1^{r+1,s}$ 在 $r=0$ 时由 $(\partial c)(b)=(-1)^{|b|}cb$ 给出，而在 $r>0$ 时由
\[
\begin{aligned}
(\partial t)(b,a_r,\ldots,a_1)
&=(-1)^{\sgnmark_r+|b|}t(b,a_r,\ldots,a_2)a_1\\
&\quad+(-1)^{|b|}t(ba_r,a_{r-1},\ldots,a_1)\\
&\quad+\sum_n(-1)^{\triangle} t(b,a_r,\ldots,a_{n+2}a_{n+1},\ldots,a_1)
\end{aligned}
\]
\SourcePage{39}
给出，其中 $\triangle=|a_{n+2}|+\cdots+|a_r|+|b|-r+n+1$。这是标准的条形分解。当 $A$ 具有恒等态射，且这些态射在 $M$ 上作为恒等映射作用时，它具有收缩同伦
\[
(\kappa t)(b,a_{r-1},\ldots,a_1)=t(e_Y,b,a_{r-1},\ldots,a_1).
\]
因此\eref{2.12}无上同调，从而得到所需结论。
\end{proof}

专门考察来自 Yoneda 嵌入的 $\Ai$ 模 $\M$，此时 $\lambda=\mathcal I^1$，因此：
\begin{corollary}\label{st:2.13}
$\mathcal I$ 上同调全忠实（从而也 c-含幺）。
\end{corollary}
考察 $\mathcal I$ 的像，还可得到：
\begin{corollary}\label{st:2.14}
任意 c-含幺 $\Ai$ 范畴，都以典范方式拟同构于一个严格含幺 dg 范畴。
\end{corollary}

现在考察 Yoneda 嵌入对于 c-含幺 $\Ai$ 范畴之间的上同调全忠实 $\Ai$ 函子 $\F:\A\to\B$ 的自然性。对任意这样的函子 $\F$，由\stref{引理}{1.16}，\eref{1.23}中自然变换 $T$ 的一阶项是 c-含幺 $\Ai$ 模的同构。结合\stref{引理}{1.6}，可知 $T$ 本身在 $\Ai$ 函子的基础上同调范畴中是同构。换言之，图
\SourceEquationDiagram{diagram-2-13}{0.28246124}{2.13}
在 $H^0(\operatorname{fun}(\A,\modcat(\A)))$ 中在典范同构意义下交换。

\begin{note}\label{st:2.15}
仿照\stref{笔记}{2.7}的思路，可以把\stref{引理}{2.12}的证明重新表述得更抽象一些。设 $\operatorname{Mod}^0(A)$ 为从 $A$ 到分次向量空间的反变（含幺）函子组成的Abel 范畴。若 $\M_0,\M_1$ 为 c-含幺 $\Ai$ 模，它们的基础上同调模 $M_0,M_1$ 就是 $\operatorname{Mod}^0(A)$ 的对象。借助条形分解，由 $\hom_{\Q}(\M_0,\M_1)$ 的长度滤过得到的谱序列，其 $E_2$ 项可解释为
\begin{equation}\label{eq:2.14}
E_2=\operatorname{Ext}_{\operatorname{Mod}^0(A)}(M_0,M_1),
\end{equation}
右侧的双分次方式与\eref{2.8}相同。与 Yoneda $\Ai$ 模关联的上同调模，是 $\operatorname{Mod}^0(A)$ 的投射对象。因此，若 $\M_0=\Y$，则\eref{2.14}只有一列非零，而这一列可与 $M_1(Y)$ 等同。所以 $\M_1(Y)$ 拟同构于 $\hom_{\Q}(\Y,\M_1)$。
\end{note}

\SourcePage{40}
\subsection{术语约定}\label{sec:2h}
关于恒等态射的问题已经讨论完毕；今后我们将不再在术语中强调这一点，而采用如下约定：\emph{除非另作说明，以下所有 $\Ai$ 范畴、$\Ai$ 函子和 $\Ai$ 模均假定为 c-含幺。}
